\section{04 Linux Kernel}

\subsection{Kernel Characteristics \& History}

\begin{concept}{Linux Kernel -- Characteristics}
    \begin{itemize}
        \item An \important{executable binary}
        \item Drivers \important{can be part} of the kernel
        \item Must be \important{recompiled} if a driver is added as built-in (\texttt{*})
        \item Linux is \important{both modular and monolithic}
    \end{itemize}
\end{concept}

\begin{example2}{Exam Quiz: What is true about the Linux kernel? (slide 2)}
    Select one or more:
    \begin{enumerate}
        \item System calls are the only way to communicate with the kernel
        \item Kernel modules run slower than built-in drivers
        \item The kernel must be recompiled when a driver is added with \texttt{"*"}
        \item Drivers can be part of the kernel
        \item It is an executable binary
    \end{enumerate}
    \tcblower
    \important{Correct: (3), (4), (5).}
    \begin{itemize}
        \item (1) false -- sysfs/\texttt{/proc} also communicate with the kernel
        \item (2) false -- modules run at the \emph{same} speed as built-in drivers
    \end{itemize}
\end{example2}

\begin{definition}{Birth of Linux}
    \begin{itemize}
        \item \important{Stallman / GNU}: tools (\texttt{gcc}, \texttt{bash}), libs (\texttt{glibc})
              -- but \emph{no kernel} (Hurd never finished)
        \item \important{Torvalds (1991)}: wrote a simple \important{monolithic} kernel
              "Linux", founded \texttt{kernel.org}
        \item Linux = GNU tools/libs \important{+} Torvalds' kernel
    \end{itemize}
\end{definition}

\begin{example2}{Exam Quiz: Components of the Linux project? (slide 4)}
    (vs the GNU project) -- mark the Linux-project components:
    \begin{enumerate}
        \item Free Software Foundation
        \item Monolithic Kernel
        \item Loadable Kernel Module
        \item gcc
        \item www.kernel.org
    \end{enumerate}
    \tcblower
    \important{Correct: (2), (3), (5).}
    FSF and \texttt{gcc} belong to \important{GNU} (Stallman), not the Linux project.
\end{example2}

\subsection{Kernel Architecture}

IMAGE (slide 8): kernel architecture -- apps / system libs (libc) / syscall + sysfs interface / kernel (scheduler, MM, IPC, FS, networking, drivers, modules) / arch-dependent code / hardware <<<

\mult{2}

\begin{concept}{Monolithic vs. Microkernel}
    \begin{itemize}
        \item \important{Monolithic} (Linux): everything in kernel space; easier to
              program \& debug
        \item \important{Microkernel} (GNU Hurd): only the essentials in the kernel;
              drivers/FS as user-space servers
    \end{itemize}
\end{concept}

\begin{corollary}{Microkernel Core Components}
    A microkernel keeps only:
    \begin{itemize}
        \item \important{Scheduler}
        \item \important{Interprocess Communication (IPC)}
        \item \important{Virtual Memory controller}
    \end{itemize}
    (File system \& device drivers live in user space.)
\end{corollary}

\multend

\begin{example2}{Exam Quiz: Microkernel (Hurd) components? (slide 6)}
    Select one or more: Scheduler $\cdot$ File System $\cdot$ IPC $\cdot$
    Device Drivers $\cdot$ Virtual Memory Controller
    \tcblower
    \important{Correct: Scheduler, IPC, Virtual Memory Controller.}
    File System \& Device Drivers are \emph{user-space servers} in a microkernel.
\end{example2}

\begin{definition}{Monolithic vs. Modular Kernel}
    \begin{itemize}
        \item \important{Monolithic}: single executable, all drivers built in; adding
              a driver $\to$ \important{recompile} the whole kernel
        \item \important{Modular}: monolithic + a \important{module interface}; add a
              driver as a module \important{at run-time}, no recompile
        \item Linux is \important{both} (modular monolithic kernel)
    \end{itemize}
\end{definition}

\begin{remark}
    \important{Recent:} PREEMPT\_RT (real-time) merged into mainline since Linux
    \important{6.12} (Nov 2024); \important{Rust} declared no longer experimental
    (2025), now a core part of the kernel.
\end{remark}

\subsection{Getting \& Building the Kernel}

\mult{2}

\begin{definition}{Getting the Sources}
    From \texttt{kernel.org} via \important{HTTP / GIT / RSYNC}; the "vanilla" kernel.
    Release types:
    \begin{itemize}
        \item \important{mainline} (rc), \important{stable}, \important{longterm}
        \item tarballs (full source) + patches (diffs between versions)
    \end{itemize}
\begin{lstlisting}[language=bash, style=basesmol, aboveskip=2pt, belowskip=2pt]
git clone \
git://git.kernel.org/.../linux.git
\end{lstlisting}
\end{definition}

\begin{definition}{Source Tree -- Key Directories}
    \begin{itemize}
        \item \texttt{init/} -- arch-independent boot code (\texttt{main.c} $\to$
              \texttt{start\_kernel()})
        \item \texttt{ipc/} -- System V IPC (\texttt{sem}/\texttt{msg}/\texttt{shm})
        \item \texttt{kernel/} -- core: \texttt{sched}, fork/exec/signal, modules
        \item \texttt{arch/} -- per-CPU ports (x86, arm, arm64, riscv\dots)
        \item \texttt{drivers/} -- \important{largest} part of the tree (\textasciitilde61\%)
        \item \texttt{fs/} -- filesystems; \texttt{include/} -- headers (\texttt{asm-*}, \texttt{linux})
    \end{itemize}
\end{definition}

\multend

\mult{2}

\begin{definition}{Kernel Configuration}
    Config stored in \texttt{.config} (text, \texttt{key=value}); never edit by hand:
    \begin{itemize}
        \item \texttt{make menuconfig} (text, needs \texttt{libncurses-dev})
        \item \texttt{make xconfig} (graphical)
        \item undo: \texttt{cp .config.old .config}
        \item Yocto: \texttt{bitbake -c menuconfig virtual/kernel}
    \end{itemize}
\end{definition}

\begin{definition}{Built-in vs. Module}
    \begin{itemize}
        \item \texttt{*} (\important{built-in}): compiled into the image, always
              available at boot, loaded by the bootloader
        \item \texttt{M} (\important{module}): loaded/unloaded \important{dynamically},
              stored in the filesystem (needs FS access); easy driver dev: load,
              test, unload, rebuild\dots
    \end{itemize}
\end{definition}

\multend

\begin{KR}{Build \& Install the Kernel}
    \textbf{Host:}
\begin{lstlisting}[language=bash, style=basesmol, aboveskip=2pt, belowskip=2pt]
make   # -> arch/<arch>/boot/{zImage,bzImage,Image}
make install               # as root: installs /boot/vmlinuz-<ver>,
                           # System.map-<ver>, config-<ver>; updates grub
make INSTALL_MOD_PATH=<dir>/ modules_install  # /lib/modules/<ver>/
\end{lstlisting}
    \textbf{Cross-compile (embedded):} pass \texttt{ARCH} (= subdir in \texttt{arch/})
    and \texttt{CROSS\_COMPILE} (tool prefix):
\begin{lstlisting}[language=bash, style=basesmol, aboveskip=2pt, belowskip=2pt]
make ARCH=arm CROSS_COMPILE=arm-linux-
\end{lstlisting}
    \textbf{Yocto} handles all compiler switches:
\begin{lstlisting}[language=bash, style=basesmol, aboveskip=2pt, belowskip=2pt]
bitbake virtual/kernel -c menuconfig   # configure
bitbake virtual/kernel                 # compile
\end{lstlisting}
\end{KR}

\subsection{Booting the Kernel}

\begin{KR}{Boot with U-Boot}
    \begin{enumerate}
        \item load \texttt{zImage} at address X
        \item load \texttt{<board>.dtb} at address Y
        \item start: \texttt{bootz X - Y} (ARM32) or \texttt{booti X - Y} (ARM64/RISC-V)
    \end{enumerate}
    U-Boot passes the \important{DTB} alongside the kernel. Kernel command line:
\begin{lstlisting}[language=bash, style=basesmol, aboveskip=2pt, belowskip=2pt]
console=ttyS0 root=/dev/mmcblk0p2 rootwait
\end{lstlisting}
    (\texttt{root=} root FS, \texttt{console=} destination of kernel messages.)
\end{KR}

\mult{2}

\begin{definition}{First Process: /sbin/init}
    First process, \important{ancestor of all}; controls runlevels:
    \begin{itemize}
        \item 0 shutdown, 1 single-user
        \item 2 multi-user (no NFS), 3 full multi-user
        \item 5 X11, 6 reboot
    \end{itemize}
    Runs startup/shutdown scripts per runlevel.
\end{definition}

\begin{definition}{Shutdown}
    Use \texttt{/bin/shutdown} -- never pull the plug (buffered writes $\to$ corrupt
    FS). \texttt{init} sends \important{SIGTERM} then \important{SIGKILL}.
    \begin{itemize}
        \item \texttt{-h} halt, \texttt{-r} reboot, \texttt{-p} poweroff
        \item Ctrl-Alt-Del defined in \texttt{/etc/inittab}
    \end{itemize}
\end{definition}

\multend

\subsection{System Calls}

IMAGE (slide 34): ARM ISR/exception vector table <<<

\begin{concept}{System Calls}
    The main interface between \important{user space} and the \important{kernel}
    (\textasciitilde300 calls). They:
    \begin{itemize}
        \item make programming easier (kernel handles HW specifics)
        \item increase security (kernel checks each request)
        \item provide portability (stable interface, changing implementation)
    \end{itemize}
    Wrapped by the C library -- \important{no direct access}.
\end{concept}

\begin{definition}{Kernel Mode vs. User Mode}
    \begin{itemize}
        \item \important{Kernel mode} (privileged): full, unrestricted HW access;
              any instruction/address; a crash \important{crashes the system}
        \item \important{User mode}: no direct HW/memory access (only via syscalls);
              crashes are \important{recoverable}; most code runs here
    \end{itemize}
\end{definition}

\begin{example2}{Exam Quiz: Kernel Mode (privileged)? (slide 33)}
    Select one or more:
    \begin{enumerate}
        \item Allows access to all underlying hardware
        \item The system can recover from program crashes
        \item Reserved for the most safety-critical OS functions
        \item Programs mostly run in privileged mode
    \end{enumerate}
    \tcblower
    \important{Correct: (1) and (3).}
    (2) false -- a kernel-mode crash takes down the system; (4) false -- most code
    runs in \emph{user} mode.
\end{example2}

\begin{KR}{Invoke a System Call on ARM}
    Triggered by a \important{supervisor call} (software interrupt) \texttt{SVC \#0}:
    \begin{itemize}
        \item \texttt{r7} = system call number (from \texttt{arch/arm/kernel/calls.S})
        \item \texttt{r0--r2} = arguments
        \item \texttt{SVC \#0} $\to$ switch to kernel mode, execute, return
    \end{itemize}
    Example -- write "Hello World" to stdout (\texttt{sys\_write} = 4):
\begin{lstlisting}[style=basesmol, aboveskip=2pt, belowskip=2pt]
HelloWorldString: .ascii "Hello World\n"
main:
mov r7, #4          /* sys_write */
mov r0, #1          /* stdout = 1 */
ldr r1, =HelloWorldString
mov r2, #12         /* string length */
svc #0
\end{lstlisting}
\end{KR}

\begin{remark2}{Modifying the Syscall Table}
    One can in principle replace an entry in \texttt{sys\_call\_table[\_\_NR\_open]}
    with a custom function (e.g. spy code). \important{But} in modern kernels the
    table is in \texttt{.rodata} (read-only) -- writing to it \important{instantly
    crashes} the system.
\end{remark2}

\raggedcolumns
\columnbreak


\subsection{eBPF (extended Berkeley Packet Filter)}

IMAGE (slide 41): eBPF flow -- BPF program -> LLVM/Clang -> bytecode -> eBPF verifier + JIT -> kernel; BPF maps shared with user space <<<

\begin{concept}{eBPF}
    Runs \important{sandboxed} programs in kernel space \important{without} kernel
    modules or source changes.
    \begin{itemize}
        \item \important{In-kernel verifier}: static analysis $\to$ can't crash/hang the kernel
        \item \important{event-driven}: attach to syscalls, packets, kprobes
        \item \important{JIT compiled} (near-native), uses LLVM; \important{no loops}
        \item user $\leftrightarrow$ kernel via eBPF \important{"maps"}
    \end{itemize}
\end{concept}

\begin{examplecode}{eBPF with BCC in Python (slide 42)}
\begin{lstlisting}[language=Python, style=basesmol, aboveskip=2pt, belowskip=2pt]
from bcc import BPF
BPF_SOURCE_CODE = r"""
TRACEPOINT_PROBE(syscalls, sys_enter_mkdirat) {
bpf_trace_printk("Wow, new directory: %s\n", args->pathname);
return 0;
}
"""
bpf = BPF(text=BPF_SOURCE_CODE)
while True:                       # event loop
(task, pid, cpu, flags, ts, msg) = bpf.trace_fields()
print("%-6d %s" % (pid, msg))
\end{lstlisting}
\end{examplecode}

\subsection{System Logging}

\begin{definition}{Kernel \& System Messages}
    \begin{itemize}
        \item \texttt{dmesg} -- kernel message \important{ring buffer} (last \textasciitilde200 msgs)
        \item \texttt{cat /var/log/messages} -- kernel + system info via the syslog daemon
    \end{itemize}
\end{definition}

\begin{formula}{Console Log Levels}
    \begin{tabular}{l|c|l}
        \texttt{KERN\_EMERG}   & 0 & system unstable / about to crash \\
        \texttt{KERN\_ALERT}   & 1 & act immediately \\
        \texttt{KERN\_CRIT}    & 2 & critical (HW/SW failure) \\
        \texttt{KERN\_ERR}     & 3 & error (often driver HW issues) \\
        \texttt{KERN\_WARNING} & 4 & warning \\
        \texttt{KERN\_NOTICE}  & 5 & notable (e.g. security) \\
        \texttt{KERN\_INFO}    & 6 & informational \\
        \texttt{KERN\_DEBUG}   & 7 & debug \\
    \end{tabular}
    \vspace{1mm}\\
    Aliases in C: \texttt{pr\_emerg}\dots\texttt{pr\_debug}. View/set console level:
\begin{lstlisting}[language=bash, style=basesmol, aboveskip=2pt, belowskip=2pt]
cat /proc/sys/kernel/printk   # current default minimum boot-default
echo 8 > /proc/sys/kernel/printk   # show all logs
\end{lstlisting}
\end{formula}

\raggedcolumns
